\section{Normal modes: from coupled oscillators to the return covariance}
\label{sec:theory}

\subsection{Coupled oscillators}

Consider $N$ unit masses with displacements $x\in\mathbb{R}^N$ coupled by
springs with a symmetric positive-definite stiffness matrix $K$,
\begin{equation}
  \ddot{x} = -Kx .
  \label{eq:oscillator}
\end{equation}
Writing $K = U\Omega^2U^\top$ with $U$ orthogonal and
$\Omega^2 = \operatorname{diag}(\omega_1^2,\dots,\omega_N^2)$, the normal
coordinates $\xi = U^\top x$ decouple, $\ddot\xi_k = -\omega_k^2\xi_k$,
and every motion of the system is a superposition of $N$ independent
oscillations \citep{goldstein2002}. Two features of this textbook result
carry over to markets: the decomposition is \emph{orthogonal}, so
normal coordinates can be traded as separate portfolios, and each mode
has its own \emph{time scale}, $2\pi/\omega_k$.

\subsection{The overdamped analogue: an Ornstein--Uhlenbeck market}

Prices are not inertial, so the relevant limit of
\eqref{eq:oscillator} is the overdamped one with thermal forcing. Let
$x_t$ denote a vector of log-price deviations from some slowly moving
equilibrium (for instance, idiosyncratic mispricings relative to a
factor model) and suppose
\begin{equation}
  \mathrm{d}x_t = -K x_t\,\mathrm{d}t + s\,\mathrm{d}W_t ,
  \label{eq:ou}
\end{equation}
a multivariate Ornstein--Uhlenbeck process \citep{uhlenbeck1930} with
isotropic noise of intensity $s^2$. The normal coordinates of $K$ are
again independent, now as scalar OU processes,
\begin{equation}
  \mathrm{d}\xi_{k,t} = -\omega_k^2\,\xi_{k,t}\,\mathrm{d}t + s\,\mathrm{d}W_{k,t},
\end{equation}
each relaxing toward equilibrium at rate $\kappa_k=\omega_k^2$ (half-life
$\ln 2/\omega_k^2$). The stationary covariance $\Sigma$ solves the
Lyapunov equation $K\Sigma + \Sigma K = s^2 I$, which gives
\begin{equation}
  \Sigma = \frac{s^2}{2}\,K^{-1}
  \quad\Longrightarrow\quad
  \Sigma \text{ and } K \text{ share eigenvectors, with }
  \lambda_k = \frac{s^2}{2\,\omega_k^2}.
  \label{eq:equipartition}
\end{equation}
Equation~\eqref{eq:equipartition} is equipartition: soft modes (small
$\omega_k$) fluctuate a lot, stiff modes fluctuate little. It also turns
the covariance matrix into a map of relaxation speeds,
\begin{equation}
  \kappa_k \;=\; \frac{s^2}{2\lambda_k},
  \qquad
  \text{half-life}_k \;=\; \frac{2\ln 2}{s^2}\,\lambda_k .
  \label{eq:prediction}
\end{equation}
In this idealized market, \textbf{the lowest-variance eigenportfolios
are the fastest to mean-revert}, and the log half-life should increase
one-for-one with $\log\lambda_k$. This is the prediction we take to the
data. It is the opposite of the usual reading of a principal-component
spectrum, in which the large eigenvalues are ``signal'' and the small
ones are noise to be discarded.

Two caveats matter. First, \eqref{eq:ou} describes \emph{levels};
we observe daily returns $r_t\approx \mathrm{d}x_t$, whose covariance
is dominated by the noise term $s^2 I$ at short horizons. The
eigenstructure of the return covariance reflects the drift term only
through the (weak) autocorrelation of returns. The model is therefore
a heuristic for where reversion should be found, not a structural
model of the return covariance. Second, real return covariances are
dominated by a market mode and a handful of sector modes whose
eigenvalues are far larger than any OU drift could produce. We
interpret these as common factors (non-stationary, $\omega_k\to 0$)
and exclude the market mode from every strategy.

\subsection{Mode portfolios}

We work with the correlation matrix of returns standardized by their
in-window volatilities $\sigma_i$. Its eigenvectors $v_k$ (columns of
$V$, ordered by descending eigenvalue $\lambda_0\ge\lambda_1\ge\dots$)
define mode returns
\begin{equation}
  r_{k,t} \;=\; \sum_i v_{ik}\,\frac{r_{i,t}}{\sigma_i},
  \qquad \operatorname{Var}(r_k)\approx\lambda_k \text{ in-sample},
  \label{eq:mode-return}
\end{equation}
which are the returns of the tradable portfolios with asset weights
$v_{ik}/\sigma_i$ \citep{avellaneda2010}. The cumulative mode ``price''
$X_{k,t}=\sum_{s\le t} r_{k,s}$ is the empirical counterpart of the
normal coordinate $\xi_k$.

\section{Separating signal from noise: Marchenko--Pastur}
\label{sec:rmt}

With $T$ observations of $N$ assets, sample eigenvalues are biased and
noisy even when the true correlation matrix is the identity. In the
limit $N,T\to\infty$ with $q=N/T$ fixed, the eigenvalues of a pure-noise
correlation matrix follow the Marchenko--Pastur law
\citep{marchenko1967},
\begin{equation}
  \rho(\lambda) = \frac{\sqrt{(\lambda_+-\lambda)(\lambda-\lambda_-)}}{2\pi q\lambda},
  \qquad
  \lambda_\pm = (1\pm\sqrt q)^2,
  \label{eq:mp}
\end{equation}
with an additional point mass of weight $1-1/q$ at zero when $q>1$.
Eigenvalues above $\lambda_+$ are inconsistent with pure noise; following
\citet{laloux1999} and \citet{plerou2002}, we call them \emph{signal
modes} and denote their number by $n_s$. In our rolling windows
($N=417$, $T=252$, $q\approx1.65$) the edge is $\lambda_+\approx 5.2$ and
the sample has rank $T-1=251$.

The MP edge is a statement about eigenvalues. For trading, the stability
of eigen\emph{vectors} matters more. Spiked-model theory
\citep{benaych2011} shows that a population eigenvalue just above the
detection threshold produces a sample eigenvector that is only weakly
aligned with its population counterpart, and below the threshold the
alignment is asymptotically zero. Read literally, this argues for
trading only modes well above $\lambda_+$. Read together with
\eqref{eq:prediction}, it creates a tension that the paper is designed
to resolve empirically. The modes that RMT certifies as real are the
\emph{soft} ones that should revert slowly. The modes that the OU
picture says should revert fastest sit at the bottom of the spectrum,
inside the noise bulk, where individual eigenvectors are unreliable. Even
so, the subspace they span (the directions of least variance in the
market) can be stable.
